\usepackage[utf8]{inputenc} % Codificação do documento
\usepackage[T1]{fontenc} % Seleção de códigos de fonte
\usepackage{lmodern} % Usa a fonte Latin Modern
\usepackage[brazil]{babel} % Traduz termos para o português
\usepackage{microtype} % Melhora a paginação/justificação do texto
\usepackage{indentfirst} % Indenta o primeiro parágrafo de cada seção

\usepackage[alf]{abntex2cite}

\pagestyle{empty}

\usepackage{multicol}

\usepackage[tmargin=1cm, bmargin=3cm, lmargin=2.5cm, rmargin=2.5cm]{geometry}
\usepackage{multicol,multirow}
\setlength{\columnsep}{.75cm}
\usepackage{xcolor}


\usepackage{tikz}
\usetikzlibrary{positioning, arrows.meta}


% ESPAÇAMENTO ENTRE PARAGRÁFOS
\usepackage{parskip}
\setlength{\parskip}{1em}

% RECUO
\setlength{\parindent}{1.25cm}

\usepackage{textcomp}

\usepackage{graphicx}

\usepackage{setspace}
\usepackage{float} 

\usepackage{array}
\usepackage{colortbl}

%\pgfplotsset{compat=1.18}

\usepackage{tikz}
\usepackage{pgfplots} % Pacote para plotagem de gráficos
\pgfplotsset{compat=1.18} % Configurações opcionais
%\usetkzobj{all}
\usepackage{adjustbox}
\usepackage{enumitem} % Pacote para personalizar listas
\usepackage{pifont} % Pacote para símbolos adicionais
\usepackage{array}

\setlength{\parindent}{0pt}

\usepackage{amssymb,amsmath,amsfonts,amsthm,mathtools,bm,siunitx,empheq,gensymb}

%\usepackage{ntheorem}

%\usepackage{fancyhdr}

\usepackage[utf8]{inputenc}
\usepackage{tikz}
\usepackage{amsmath,amssymb}
\usetikzlibrary{patterns,arrows.meta,shapes.geometric,positioning,calc,decorations.pathreplacing,shadows}

\usepackage{amsmath,amssymb}
\usetikzlibrary{angles,quotes,calc,arrows.meta}

\usepackage[utf8]{inputenc}
\usepackage{amsmath}
\usetikzlibrary{patterns, patterns.meta, calc, arrows.meta}

\usetikzlibrary{angles, quotes} 

\everymath{\displaystyle}

\newcommand{\aula}[4]{%
  \noindent \textbf{Aula #1} --- #2 de #3 de #4%
}

% 1. Remove a quebra de página automática dos capítulos [2]
\renewcommand{\clearforchapter}{\vspace{1.5cm}}

% 2. Opcional: reduz o espaço em branco grande acima do título do capítulo [3]
\setlength{\beforechapskip}{1em}

\newcommand{\paracasa}{%
    \vspace{16pt} %
    \textcolor{blue!80!black}{Continuação - para casa} %
}

\usetikzlibrary{patterns.meta}

\usepackage{cancel}

\usepackage{environ}

% Chave liga/desliga (true para mostrar respostas, false para ocultar)
\newif\ifmostrarrespostas
\mostrarrespostastrue % <- Altere para \mostrarrespostasfalse para gerar a versão do aluno!

\NewEnviron{resposta}{%
  \ifmostrarrespostas
    %\par\addvspace{0.5em}
    \color{gray!80!black}
    \noindent \BODY
    %\par\addvspace{0.5em}
  \fi
}

\NewEnviron{duvida}{%
  \ifmostrarrespostas
    %\par\addvspace{0.5em}
    \color{red!80!black}
    \noindent \BODY
    %\par\addvspace{0.5em}
  \fi
}

\NewEnviron{observacao}{%
  \ifmostrarrespostas
    %\par\addvspace{0.5em}
    \color{orange!80!black}
    \noindent \BODY
    %\par\addvspace{0.5em}
  \fi
}

\endinput